\section{Example: User Interfaces} \label{ex77}

Although \W offers a variety of options, sometimes it's necessary to
use some special input data, to write a user-defined output, to perform
a complicated fit, or to provide some unusual field map, or {\ldots}

There are some routines to be provided by the user:

\begin{itemize}
\item \gti{subroutine uname}

This routine allows to provide user input stuff or access and modify
the variables of \We.

\item \gti{
subroutine uout}

You can use this routine to access the variables of WAVE at the end of
the run and to use it for your own output.

\item \gti{
subroutine ustep(x,y,z,vx,vy,vz,vxp,vyp,vzp,dt,gamma,dgamma)}

This routine is called during each tracking step. You can manipulate
the given arguments.

\item \gti{
subroutine bextern(x,y,z,bx,by,bz,ax,ay,az)}

If you want to provide your own magnetic field, set \scap(KBEXTERN=1).
The routine must return the magnetic field and vector potential (if
needed) for the point \gti{(x,y,z)}.

\item \gti{
subroutine ubmask(x,y,z,bxmask,bymask,bzmask)}

This routine multiplies the magnetic field (\scap{IBMASK=-100}) or
adds the masking field (\scap{IBMASK=-200}). Please check the namelist
\scap{\$B0SCBLOBN} in \wie.

\end{itemize}


